% Part: incompleteness
% Chapter: introduction
% Section: historical-background

\documentclass[../../../include/open-logic-section]{subfiles}

\begin{document}

\olfileid[hi]{inc}{int}{bgr}

\olsection{ऐतिहासिक पृष्ठभूमि}

इस अनुभाग में हम उन ऐतिहासिक विकासों की संक्षेप में चर्चा करेंगे जिनसे
अपूर्णता प्रमेयों को उनके संदर्भ में समझने में सहायता मिलेगी। विशेष रूप से,
हम गणितीय तर्कशास्त्र के इतिहास का बहुत स्थूल अवलोकन देंगे और फिर गणित की
आधारशिलाओं के इतिहास के विषय में कुछ बातें कहेंगे।

\begin{digress}
  ``गणितीय तर्कशास्त्र'' पद द्व्यर्थी है। ``गणितीय'' शब्द को विषय-वस्तु का
  वर्णन मानकर इसका अर्थ ``गणित का तर्कशास्त्र'', अर्थात् गणितीय तर्कणा के
  सिद्धांत, लिया जा सकता है; अथवा इसे विधियों का वर्णन मानकर ``तर्कशास्त्र
  का गणित'', अर्थात् तर्कणा के सिद्धांतों का गणितीय अध्ययन, समझा जा सकता
  है। आगे का विवरण गणितीय तर्कशास्त्र को इन दोनों अर्थों में, प्रायः एक ही
  समय, ग्रहण करता है।
\end{digress}

तर्कशास्त्र का अध्ययन मूलतः अरस्तू से आरंभ हुआ, जो लगभग 384--322
\textsc{ईसा-पूर्व} जीवित थे। उनकी \emph{कोटियाँ}, \emph{पूर्व
  विश्लेषिकी} और \emph{उत्तर विश्लेषिकी} में वैज्ञानिक तर्कणा के सिद्धांतों
का क्रमबद्ध अध्ययन मिलता है, जिसमें न्यायवाक्य का विस्तृत और सुव्यवस्थित
अध्ययन भी सम्मिलित है।

अरस्तू का तर्कशास्त्र पूरे मध्ययुग में शास्त्रीय दर्शन पर छाया रहा; वास्तव
में अठारहवीं शताब्दी के उत्तरार्ध तक कांट का मत था कि अरस्तू का तर्कशास्त्र
पूर्ण है और उसमें किसी संशोधन की आवश्यकता नहीं। किंतु न्यायवाक्य का सिद्धांत
इतना सीमित है कि वह गणितीय तर्कणा के केवल अत्यंत सतही पक्षों का ही प्रतिरूपण
कर सकता है। इससे एक शताब्दी पहले, न्यूटन के समकालीन लाइबनित्स ने तार्किक
तर्कणा के लिए एक पूर्ण ``कलन'' की कल्पना की थी और ऐसे कलन की रचना की दिशा
में कुछ आरंभिक कदम उठाए थे; उनका वर्णन मूलतः प्रतिज्ञप्तिक तर्कशास्त्र के एक
रूप के समान था।

उन्नीसवीं शताब्दी तर्कशास्त्र के लिए निर्णायक मोड़ थी। 1854 में जॉर्ज बूल ने
\emph{चिंतन के नियम} लिखी, जिसमें प्रतिज्ञप्तिक तर्कशास्त्र का विस्तृत बीजीय
अध्ययन आधुनिक प्रस्तुतियों से बहुत दूर नहीं है। 1879 में गोटलोब फ़्रेगे ने
अपनी \emph{संकल्पना-लिपि} प्रकाशित की, जो प्रतिज्ञप्तिक तर्कशास्त्र में
परिमाणक और संबंध जोड़ती है और इस प्रकार प्रथम-कोटि तर्कशास्त्र को समाहित
करती है। वास्तव में फ़्रेगे के तार्किक तंत्रों में उच्चतर-कोटि तर्कशास्त्र
और उससे भी अधिक सामग्री थी। अपनी \emph{अंकगणित के मूल नियम} में फ़्रेगे ने
यह दिखाने का प्रयास किया कि संपूर्ण अंकगणित को उनकी संकल्पना-लिपि में केवल
तार्किक मान्यताओं से व्युत्पन्न किया जा सकता है। दुर्भाग्य से ये मान्यताएँ
असंगत निकलीं, जैसा रसेल ने 1902 में दिखाया। किंतु उस असंगत अभिगृहीत को अलग
रख दें तो फ़्रेगे ने लगभग अकेले ही आधुनिक तर्कशास्त्र का आविष्कार कर दिया—यह
एक चकित कर देने वाली उपलब्धि थी। बूल के बाद बीजीय दृष्टि रखने वाले विचारकों,
जिनमें पियर्स और श्रेडर (Schr\"oder) भी थे, ने परिमाणन तर्कशास्त्र का स्वतंत्र
रूप से विकास किया।

अब हम गणित की आधारशिलाओं से जुड़े विकासों की ओर मुड़ें। निस्संदेह, गणित में
तर्कशास्त्र की महत्त्वपूर्ण भूमिका होने के कारण इनका अभी वर्णित विकासों से
गहरा अंतःसंबंध है। उदाहरण के लिए, फ़्रेगे ने अपना तर्कशास्त्र इस स्पष्ट
उद्देश्य से विकसित किया था कि समस्त गणित को केवल उनके तार्किक ढाँचे पर आधारित
दिखाया जा सके; विशेषतः वे दिखाना चाहते थे कि गणित कांट के कथनानुसार
प्रागनुभविक \emph{संश्लेषणात्मक} सत्यों का नहीं, बल्कि प्रागनुभविक
\emph{विश्लेषणात्मक} सत्यों का समुच्चय है।

बहुत-से लोग वास्तविक गणित का जन्म यूनानियों के साथ मानते हैं। लगभग 300
ईसा-पूर्व लिखी यूक्लिड की \emph{तत्त्व} अपनी कठोरता और परिशुद्धता पर बल देने
के कारण यूनानी गणित की पहले से ही परिपक्व प्रतिनिधि है। यूक्लिड की
\emph{तत्त्व} की परिभाषाएँ और प्रमाण आज के माध्यमिक-विद्यालय ज्यामिति के
पाठ्यक्रमों में लगभग अक्षुण्ण बने हुए हैं (जहाँ तक विद्यालयों में अब भी
ज्यामिति पढ़ाई जाती है)। गणितीय तर्कणा का यह प्रतिमान केवल गणित में ही नहीं,
दर्शन की शाखाओं में भी कठोर युक्ति का आदर्श माना गया है। (स्पिनोज़ा ने तो
नैतिक और धार्मिक तर्क भी यूक्लिडीय शैली में प्रस्तुत किए थे—उन्हें इस रूप में
देखना विचित्र लगता है!)

सत्रहवीं शताब्दी में न्यूटन और लाइबनित्स ने कलन का आविष्कार किया। (प्राथमिकता
को लेकर उग्र विवाद सदियों चला, किंतु आज अधिकांश विद्वान मानते हैं कि दोनों
विकास अधिकतर स्वतंत्र थे।) कलन में, उदाहरणतः, अनंतसूक्ष्म राशियों के अनंत
योगों के विषय में तर्क किया जाता है; इन विशेषताओं ने बिशप बर्कली की आलोचना को
बल दिया। उनका तर्क था कि ईश्वर में विश्वास उनके समय के गणित से कम युक्तिसंगत
नहीं था। अठारहवीं शताब्दी में कलन की विधियों का व्यापक उपयोग हुआ; उदाहरण के
लिए लियोनार्ड ऑयलर ने अनंत योगों वाले परिकलनों से नाटकीय परिणाम निकाले।

उन्नीसवीं शताब्दी में गणितज्ञों ने कलन को अधिक दृढ़ आधार देकर बर्कली की
आलोचनाओं का उत्तर देने का प्रयास किया। कॉशी, वाइयरश्ट्रास, बोल्ज़ानो और अन्य
लोगों के प्रयासों से सीमा, सातत्य, अवकलन और समाकलन की हमारी आधुनिक परिभाषाएँ
``एप्सिलॉन और डेल्टा'' के पदों में बनीं—दूसरे शब्दों में, उनमें अनंतसूक्ष्म
राशियों का कोई उल्लेख नहीं था। शताब्दी के उत्तरार्ध में गणितज्ञों ने और आगे
बढ़कर वास्तविक संख्याओं सहित कलन के सभी पक्षों को प्राकृतिक संख्याओं के
पदों में समझाने का प्रयास किया। (क्रोनेकर: ``ईश्वर ने पूर्ण संख्याएँ बनाईं;
शेष सब मनुष्य का कार्य है।'') 1872 में डेडेकिंड ने ``सातत्य और अपरिमेय
संख्याएँ'' लिखी, जिसमें उन्होंने दिखाया कि वास्तविक संख्याओं को परिमेय
संख्याओं के समुच्चयों के रूप में कैसे ``निर्मित'' किया जा सकता है (और जैसा
आप जानते हैं, परिमेय संख्याओं को प्राकृतिक संख्याओं के युग्मों के रूप में
देखा जा सकता है)। 1888 में उन्होंने ``प्राकृतिक संख्याएँ क्या हैं और उन्हें
क्या होना चाहिए?'' लिखी, जिसका उद्देश्य प्राकृतिक संख्याओं को शुद्ध
``तार्किक'' पदों में समझाना था। 1887 में क्रोनेकर ने ``\"Uber den
Zahlbegriff'' (``संख्या की संकल्पना पर'') लिखी, जिसमें उन्होंने सभी गणितीय
वस्तुओं को पूर्णांकों के पदों में निरूपित करने की बात कही; 1889 में जुज़ेप्पे
पियानो ने प्राकृतिक संख्याओं के लिए औपचारिक, सांकेतिक अभिगृहीत दिए।

उन्नीसवीं शताब्दी का अंत अनंत से निपटने में एक नया साहस भी लाया। उससे पहले
अनंतात्मक वस्तुओं और संरचनाओं (जैसे प्राकृतिक संख्याओं का समुच्चय) को बहुत
सावधानी से ग्रहण किया जाता था; ``अनंततः अनेक'' का अर्थ ``जितने आप चाहें'' और
``सीमा में अभिगमन'' का अर्थ ``जितना निकट आप चाहें'' समझा जाता था। किंतु जॉर्ग
कैंटर ने दिखाया कि अनंत को उसके अंकित मूल्य पर स्वीकार करना संभव है। कैंटर,
डेडेकिंड और अन्य लोगों के कार्य ने गणित की उस सामान्य समुच्चय-सैद्धांतिक समझ
को जन्म देने में सहायता की जो अब व्यापक रूप से स्वीकृत है।

इससे हम तर्कशास्त्र और आधारशिलाओं में बीसवीं शताब्दी के विकासों तक पहुँचते
हैं। 1902 में रसेल ने फ़्रेगे के तार्किक तंत्र में विरोधाभास खोजा। 1904 में
ज़र्मेलो ने तथाकथित ``वरण अभिगृहीत'' का उपयोग करके कैंटर का सुव्यवस्थापन
सिद्धांत सिद्ध किया; इस अभिगृहीत की वैधता पर बहुत बहस हुई। 1910 से 1913 के
बीच रसेल और व्हाइटहेड की \emph{प्रिंसिपिया मैथमेटिका} के तीन खंड प्रकाशित
हुए, जिन्होंने गणित को तार्किक आधार पर स्थापित करने के फ़्रेगे-कार्यक्रम को
आगे बढ़ाया। दुर्भाग्य से रसेल और व्हाइटहेड को दो ऐसे सिद्धांत स्वीकार करने
पड़े जिन्हें शुद्ध तार्किक मानना कठिन था: अनंतता का अभिगृहीत और
``अपचयनीयता'' का अभिगृहीत। 1900 के दशक में प्वाँकारे (Poincar\'e) ने गणित में
``अपूर्वसूचक परिभाषाओं'' के उपयोग की आलोचना की; और 1910 के दशक में ब्राउवर ने
संपूर्ण गणित को ऐसे ``अंतःप्रज्ञावादी'' आधार पर पुनःस्थापित करने का प्रस्ताव
देना आरंभ किया जिसमें मध्यनिषेध नियम ($!A \lor \lnot !A$) का उपयोग नहीं होता।

वे सचमुच विचित्र दिन थे!{} समस्त गणित को तर्कशास्त्र में अपचयित करने का
कार्यक्रम अब ``तर्कवाद'' कहलाता है और ऊपर बताई कठिनाइयों के कारण सामान्यतः
असफल माना जाता है। अंतःप्रज्ञावादी मानसिक रचनाओं के पदों में गणित विकसित करने
का कार्यक्रम ``अंतःप्रज्ञावाद'' कहलाता है और माना जाता है कि वह दैनिक गणित पर
अत्यधिक कठोर प्रतिबंध लगाता है। शताब्दी के मोड़ के आसपास सर्वकालीन सबसे
प्रभावशाली गणितज्ञों में से एक, डेविड हिल्बर्ट, कैंटर और डेडेकिंड द्वारा लाई
गई नई अमूर्त विधियों के प्रबल समर्थक थे: ``कैंटर ने हमारे लिए जो स्वर्ग रचा
है, उससे कोई हमें बाहर नहीं निकाल सकेगा।'' साथ ही वे इन नई विधियों की आधारगत
आलोचनाओं के प्रति संवेदनशील थे (विचित्र रूप से इन्हीं विधियों को अब
``चिरसम्मत'' कहा जाता है)। उन्होंने दोनों लाभ एक साथ पाने का यह उपाय सुझाया:
\begin{enumerate}
\item चिरसम्मत विधियों को औपचारिक अभिगृहीतों और नियमों से निरूपित करें;
  गणितीय प्रश्नों को किसी अभिगृहीतीय तंत्र में !!{formula} के रूप में निरूपित
  करें।
\item सुरक्षित, ``परिमितवादी'' विधियों का उपयोग करके सिद्ध करें कि ये
  औपचारिक निगमनात्मक तंत्र संगत हैं।
\end{enumerate}

हिल्बर्ट के कार्य ने पहले लक्ष्य की सिद्धि की दिशा में बहुत प्रगति की। 1899
में उन्होंने अपनी प्रसिद्ध पुस्तक \emph{ज्यामिति की आधारशिलाएँ} में ज्यामिति
के लिए यह कार्य कर दिया था। बाद के वर्षों में उन्होंने और उनके अनेक छात्रों
तथा सहयोगियों ने गणित के अन्य क्षेत्रों में वही करने का प्रयास किया जो
हिल्बर्ट ने ज्यामिति के लिए किया था। हिल्बर्ट ने स्वयं अंकगणित और विश्लेषण के
लिए अभिगृहीत-तंत्र दिए। ज़र्मेलो ने समुच्चय सिद्धांत का अभिगृहीतीकरण दिया,
जिसे फ्रेंकेल, स्कोलेम, फ़ॉन न्यूमान और अन्य लोगों ने विस्तृत किया। 1920 के
दशक के मध्य तक ``समस्त'' गणित के अभिगृहीतीकरण की पदवी के दो दावेदार थे:
रसेल और व्हाइटहेड की \emph{प्रिंसिपिया मैथमेटिका}, और वह सिद्धांत जो
ज़र्मेलो--फ्रेंकेल समुच्चय सिद्धांत कहलाया।

1921 में हिल्बर्ट ने इन तंत्रों की संगतता सिद्ध करने के लक्ष्य पर एक शोध
परियोजना आरंभ की। उनके कई छात्रों ने इसमें सहायता की, विशेषतः बर्नाइस,
आकरमान और बाद में गेंट्सेन ने। लक्ष्य पूरा करने का मूल विचार यह था कि गणित
में किसी असंगतता की !!{derivation} की संभावना के प्रश्न को चिह्नों के संभावित
अनुक्रमों की एक संयोजकीय समस्या में बदला जाए। अर्थात् वाक्यों के ऐसे संभावित
अनुक्रमों पर विचार किया जाए जो अंकगणित, विश्लेषण या समुच्चय सिद्धांत के किसी
अभिगृहीत-तंत्र के अभिगृहीतों से, मान लें, $!A \land \lnot !A$ की सही
!!{derivation} होने की कसौटी पूरी करते हों। चिह्नों के ऐसे अनुक्रम की
असंभाव्यता का प्रमाण—चूँकि वह स्वयं गणितीय प्रमाण है—इन अभिगृहीतीय तंत्रों में
औपचारिक बनाया जा सकता था। दूसरे शब्दों में, कोई वाक्य $\OCon$ होता जो, मान
लें, अंकगणित की संगतता कथित करता। इसके अतिरिक्त यह वाक्य संबंधित तंत्रों में
सिद्ध होना चाहिए था, विशेषतः तब जब उसके प्रमाण को केवल बहुत सीमित,
``परिमितवादी'' साधनों की आवश्यकता हो।

दूसरे लक्ष्य—कि विकसित अभिगृहीत-तंत्र प्रत्येक गणितीय प्रश्न का निपटारा
करेंगे—को दो प्रकार से सटीक बनाया जा सकता है। पहले प्रकार में इसे ऐसे सूत्रित
कर सकते हैं: गणित के किसी अभिगृहीत-तंत्र की भाषा के प्रत्येक वाक्य~$!A$ के
लिए या तो $!A$ अथवा $\lnot !A$ अभिगृहीतों से सिद्ध है। यदि यह सत्य होता, तो
ऐसा कोई वाक्य न होता जिसे अभिगृहीतों के आधार पर न सिद्ध किया जा सके, न
खंडित; ऐसा कोई प्रश्न न होता जिसका अभिगृहीत निपटारा न करते। इस गुण वाले
अभिगृहीत-तंत्र को \emph{पूर्ण} कहते हैं। निस्संदेह, किसी दिए हुए वाक्य के लिए
यह निर्धारित करना फिर भी कठिन हो सकता था कि दोनों विकल्पों में कौन-सा लागू
होता है। किंतु सिद्धांततः ऐसा करने की कोई विधि होनी चाहिए थी। वास्तव में
हिल्बर्ट द्वारा विचारित अभिगृहीत और !!{derivation} तंत्रों के लिए पूर्णता से
ऐसी विधि का अस्तित्व फलित होता—यद्यपि हिल्बर्ट को इसका ज्ञान नहीं था। प्रश्न
की दूसरी व्याख्या यह अधिक प्रबल अपेक्षा होगी: कोई यांत्रिक, संगणकीय विधि हो
जो दिए गए वाक्य~$!A$ के लिए निर्धारित करे कि वह अभिगृहीतों से व्युत्पन्न है
या नहीं।

1931 में ग्योडेल (G\"odel) ने दो ``अपूर्णता प्रमेय'' सिद्ध किए, जिनसे पता चला
कि यह कार्यक्रम सफल नहीं हो सकता। गणित का कोई पूर्ण अभिगृहीत-तंत्र नहीं है;
विशेषतः अभिगृहीतों की संगतता व्यक्त करने वाला वाक्य न सिद्ध किया जा सकता है,
न खंडित।

इससे हिल्बर्ट के मूल कार्यक्रम को घातक आघात लगा। किंतु, जैसा गणित में प्रायः
होता है, इससे शोध के रोमांचक नए मार्ग भी खुले। यदि गणित का कोई एक सर्वग्राही
औपचारिक तंत्र नहीं है, तो अधिक परिसीमित तंत्र विकसित करना और यह जाँचना उचित
है कि उनमें क्या सिद्ध किया जा सकता है। उन तंत्रों की संगतता स्थापित करने के
लिए कम प्रतिबंधित प्रमाण-विधियाँ विकसित करना और उनकी संगतता सिद्ध करने की
कठिनाई मापने के उपाय खोजना भी सार्थक है। चूँकि ग्योडेल (G\"odel) ने दिखाया
कि (लगभग) प्रत्येक औपचारिक तंत्र में ऐसे प्रश्न हैं जिनका वह निपटारा नहीं कर
सकता, इसलिए किसी दिए हुए औपचारिक तंत्र के लिए ऐसे ``रोचक'' प्रश्न खोजना और
यह निर्धारित करना सार्थक है कि उन्हें निपटाने के लिए तंत्र कितना प्रबल होना
चाहिए। आज तक तर्कशास्त्री एक नए गणितीय अनुशासन—प्रमाण सिद्धांत—में इन्हीं
प्रश्नों का अनुसरण कर रहे हैं।

\end{document}
